\chapter{详细设计}

\section{网络地址设计}

地址规划应保证结构清晰、便于扩展，并避免不同业务区域之间出现地址冲突。本文采用私有地址段进行统一规划，通过网关地址、虚拟局域网编号和业务名称建立对应关系，便于后期维护。

\section{路由与冗余设计}

核心层与汇聚层之间应配置冗余链路，避免单点故障导致大范围网络中断。核心设备之间可以采用链路聚合或动态路由协议提升可靠性。路由设计的目标是在保证可达性的基础上减少不必要的广播和无效转发。

\begin{equation}
  A = \frac{T_{\mathrm{normal}}}{T_{\mathrm{total}}}
\end{equation}

式中，$A$ 表示网络可用性，$T_{\mathrm{normal}}$ 表示网络正常运行时间，$T_{\mathrm{total}}$ 表示统计周期总时间。

\section{安全策略设计}

校园网安全策略应覆盖接入控制、区域隔离、边界防护和日志审计。对于办公区和服务器区，应限制来自宿舍区和公共无线区的直接访问。对于必须开放的服务，应通过最小权限原则设置访问控制列表。

\subsection{访问控制设计}

访问控制列表可以根据源地址、目的地址、协议类型和端口号限制访问行为。例如，服务器区只允许教学区和办公区访问特定业务端口，而禁止普通终端直接访问数据库端口。

\subsection{运维管理设计}

运维管理应包括设备配置备份、故障告警、链路状态监测和日志留存。对于关键设备，应建立定期巡检制度，并通过统一命名规则降低维护难度。

